% =====================================================================
%  Comparison of MuscleMimic, mjlab, and Qflex
%  Compile with:  pdflatex framework_comparison.tex   (run twice for refs)
%  Uses only standard CTAN packages; no shell-escape required.
% =====================================================================
\documentclass[11pt,a4paper]{article}

\usepackage[T1]{fontenc}
\usepackage[utf8]{inputenc}
\usepackage{lmodern}
\usepackage[margin=2.4cm]{geometry}
\usepackage{amsmath,amssymb}
\usepackage{booktabs}
\usepackage{array}
\usepackage{ragged2e}
\usepackage{enumitem}
\usepackage{xcolor}
\usepackage{listings}
\usepackage{caption}
\usepackage{fancyhdr}
\usepackage[hidelinks]{hyperref}

% ---------- colours ----------
\definecolor{mmC}{HTML}{AE332B}   % MuscleMimic
\definecolor{mjC}{HTML}{1B6B77}   % mjlab
\definecolor{qfC}{HTML}{8A6524}   % Qflex
\definecolor{rulegrey}{HTML}{888888}
\definecolor{codebg}{HTML}{F5F5F5}

% ---------- listings ----------
\lstset{
  basicstyle=\ttfamily\footnotesize,
  backgroundcolor=\color{codebg},
  frame=single,
  rulecolor=\color{rulegrey},
  breaklines=true,
  columns=fullflexible,
  keepspaces=true,
  showstringspaces=false,
  xleftmargin=0pt,
  framexleftmargin=3pt,
}

% ---------- helpers ----------
\newcommand{\MM}{\textcolor{mmC}{\textbf{MuscleMimic}}}
\newcommand{\MJ}{\textcolor{mjC}{\textbf{mjlab}}}
\newcommand{\QF}{\textcolor{qfC}{\textbf{Qflex}}}
\newcommand{\yes}{\textbf{yes}}
\newcommand{\no}{no}
\newcommand{\na}{\textemdash}
\newcolumntype{L}[1]{>{\RaggedRight\arraybackslash}p{#1}}

\captionsetup{font=small,labelfont=bf}

\pagestyle{fancy}
\fancyhf{}
\rhead{\small MuscleMimic / mjlab / Qflex}
\lhead{\small Framework comparison}
\cfoot{\thepage}
\renewcommand{\headrulewidth}{0.4pt}

% =====================================================================
\begin{document}

\title{\vspace{-1.5cm}\textbf{Three Roads to Musculoskeletal Control}\\[0.3em]
\large A technical comparison of MuscleMimic, mjlab, and Qflex}
\author{claude}
\date{\today}
\maketitle
\vspace{-1cm}

\begin{abstract}
\noindent
These three codebases all learn continuous control for high-dimensional, over-actuated
systems in MuJoCo. Two of them --- MuscleMimic and Qflex --- have been demonstrated on
muscle-driven human models; mjlab implements muscle actuators but ships no biomechanical
assets and has published no musculoskeletal result of its own.
They disagree at almost every design decision that matters: the learning
paradigm (on-policy policy-gradient versus off-policy value-based), the exploration
mechanism (isotropic Gaussian versus value-guided flow), the task specification
(large motion-capture corpora versus single reference trajectories versus pure task
reward), the physics backend (GPU Warp kernels versus CPU MuJoCo with process-level
parallelism), the numerics framework (JAX/Flax, PyTorch, JAX/Haiku), and the artifact
produced at the end (a sharded training checkpoint, a portable ONNX graph, a Python
pickle). A separate section covers \textbf{MyoSuite}, which is not a fourth competitor but
the substrate two of the three build on, and whose deliberately algorithm-agnostic design
explains much of the divergence above. This report documents those differences, states what
was verified directly against the source, and closes with the practical consequences for
policy distillation and hardware deployment.
\end{abstract}

\tableofcontents
\vspace{0.6em}
\hrule
\vspace{0.8em}

% =====================================================================
\section{Scope and provenance}

\begin{table}[h!]
\centering
\small
\begin{tabular}{@{}lL{3.3cm}L{3.4cm}L{3.6cm}@{}}
\toprule
 & \textcolor{mmC}{\textbf{MuscleMimic}} & \textcolor{mjC}{\textbf{mjlab}} & \textcolor{qfC}{\textbf{Qflex}} \\
\midrule
Origin      & EPFL (Mathis Lab) & MuJoCo Lab / community & Tsinghua (Sui Lab) \\
Publication & arXiv 2603.25544 (2026) & no paper; v1.4.0 release & ICLR 2026, arXiv 2601.19707 \\
Package     & \texttt{musclemimic} + vendored \texttt{loco\_mujoco} & \texttt{mjlab} & \texttt{relax} \\
Lineage     & fork of LocoMuJoCo v1.0.1 & Isaac Lab API on MuJoCo Warp & standalone \\
What it is  & an \emph{application}: MSK motion-imitation at scale & a \emph{framework}: general RL/robotics environment toolkit & an \emph{algorithm}: a new exploration mechanism, plus a benchmark harness \\
\bottomrule
\end{tabular}
\caption{What each project actually is. The category difference drives most of what follows:
MuscleMimic optimises one task family deeply, mjlab provides reusable machinery,
and Qflex contributes a learning rule evaluated across borrowed benchmarks.}
\end{table}

\noindent
This asymmetry is worth holding onto. Comparing them like-for-like is slightly unfair:
MuscleMimic ships two validated musculoskeletal models and a retargeted corpus but only
one algorithm; mjlab ships a rich environment API and deployment tooling but no
biomechanical assets; Qflex ships twelve algorithms and no environments of its own
(every benchmark is a git submodule).

% =====================================================================
\section{Executive summary}

\begin{table}[h!]
\centering
\footnotesize
\begin{tabular}{@{}lL{3.2cm}L{3.2cm}L{3.5cm}@{}}
\toprule
 & \textcolor{mmC}{\textbf{MuscleMimic}} & \textcolor{mjC}{\textbf{mjlab}} & \textcolor{qfC}{\textbf{Qflex}} \\
\midrule
Learning paradigm & on-policy policy gradient & on-policy policy gradient & \textbf{off-policy actor--critic} \\
Algorithm & PPO (vendored, PureJaxRL lineage) & PPO (\texttt{rsl\_rl} dependency) & \textbf{Q-guided flow exploration} \\
Exploration & Gaussian, learnable $\sigma$ & Gaussian, learnable $\sigma$ & \textbf{value-gradient flow} \\
Data reuse & 1 gradient epoch, discard & 5 epochs typical, discard & \textbf{replay buffer, $10^6$} \\
Task signal & 972+ mocap clips & command or 1 motion & task reward (+1 ref.\ traj) \\
Physics & MuJoCo Warp via MJX & MuJoCo Warp direct & \textbf{CPU MuJoCo} \\
Numerics & JAX + Flax & \textbf{PyTorch} & JAX + Haiku \\
Parallelism & 8192 GPU envs & 4096 GPU envs & \textbf{8 CPU processes} \\
Budget & $10^9$ steps ($\approx$20\,h) & $\sim\!10^8$--$10^9$ steps & $5\times10^6$ steps \\
Output & Orbax checkpoint dir & \textbf{.pt + .onnx} & dill pickle \\
Deployable artifact & \no & \yes & \no \\
Distillation & \no & available, unwired & \no \\
\bottomrule
\end{tabular}
\caption{At a glance. Bold marks the entry that differs from the other two.}
\end{table}

\noindent
Three headline observations:

\begin{enumerate}[leftmargin=1.4em,itemsep=0.35em]
\item \textbf{Qflex is the theoretical outlier.} MuscleMimic and mjlab are both on-policy
PPO with Gaussian exploration and differ mainly in engineering. Qflex changes the learning
rule itself, and its contribution is specifically an \emph{exploration} mechanism for
over-actuated systems --- exactly the regime of a 400--700 muscle model.

\item \textbf{mjlab is the engineering outlier.} It is the only one that emits a portable,
self-describing policy artifact, and the only one whose observation pipeline models
sensor latency.

\item \textbf{MuscleMimic is the data outlier.} It is the only one built around a large
retargeted motion corpus with a goal-conditioned, generalist policy; the other two train
single-task or single-clip controllers.
\end{enumerate}

% =====================================================================
\section{Theoretical learning approach}

\subsection{MuscleMimic and mjlab: on-policy PPO}

Both optimise the standard clipped surrogate objective,
\begin{equation}
L^{\text{CLIP}}(\theta) = \hat{\mathbb{E}}_t\!\left[\min\!\big(r_t(\theta)\hat{A}_t,\;
\mathrm{clip}(r_t(\theta), 1-\epsilon, 1+\epsilon)\hat{A}_t\big)\right],
\qquad r_t(\theta) = \frac{\pi_\theta(a_t|s_t)}{\pi_{\theta_{\text{old}}}(a_t|s_t)},
\end{equation}
with a clipped value loss, GAE$(\gamma,\lambda)$ advantages, and a diagonal Gaussian
policy $\pi_\theta(a|s) = \mathcal{N}(\mu_\theta(s), \mathrm{diag}(\sigma^2))$ whose
$\sigma$ is state-independent and learned. Exploration is therefore isotropic noise in
the native action space --- precisely the strategy Qflex argues degrades with dimension.

The two implementations are genealogically close: mjlab calls \texttt{rsl\_rl} as a
dependency, while MuscleMimic vendors a PPO whose lineage traces to PureJaxRL through
LocoMuJoCo. Two fingerprints confirm the latter: the GAE scan carries the constant
\texttt{unroll=16} (preserved in MuscleMimic's \texttt{rl\_core/rollout\_buffer.py} with
the comment \emph{``matching PPO''}), and LocoMuJoCo retains PureJaxRL's characteristic
misspelling \texttt{\_update\_minbatch}, which MuscleMimic silently corrected.

\paragraph{MuscleMimic's one substantive algorithmic claim.}
Its paper argues that for muscle-actuated systems the number of gradient epochs per batch
must be $E=1$ --- strictly on-policy. With $E=3$ or $E=10$ the KL divergence between the
current and data-generating policy spikes above $10^{10}$, whereas $E=1$ stays below
$10^{-1}$. The attributed cause is muscle activation dynamics,
\begin{equation}
\frac{\partial\,\mathrm{act}}{\partial t} = \frac{\mathrm{ctrl}-\mathrm{act}}{\tau(\mathrm{ctrl},\mathrm{act})},
\qquad
\tau = \begin{cases}
\tau_{\text{act}}(0.5 + 1.5\,\mathrm{act}), & \mathrm{ctrl} > \mathrm{act}\\
\tau_{\text{deact}}/(0.5 + 1.5\,\mathrm{act}), & \mathrm{ctrl} \le \mathrm{act}
\end{cases}
\end{equation}
which amplifies small policy changes through a first-order lag
($\tau_{\text{act}}=10$\,ms, $\tau_{\text{deact}}=40$\,ms). Discarding every sample after
one gradient step is affordable only because GPU simulation is cheap; the compensating
lever is $N_{\text{env}}$, not $E$.

\subsection{Qflex: off-policy actor--critic with value-guided flow}

Qflex replaces Gaussian exploration with a learned transport. The motivating observation
is that isotropic perturbations lose coverage as $|\mathcal{A}|$ grows, so exploration
should follow task-relevant gradients instead. It constructs a velocity field aligned
with the action-gradient of the critic,
\begin{equation}
\frac{da^{(t)}}{dt} = v_Q^{(t)}(a^{(t)};s) = M\,\nabla_a Q^{\pi_{\text{old}}}(s,a),
\qquad a^{(0)} \sim \pi^{(0)}(\cdot|s),
\end{equation}
with $M$ a positive-definite preconditioner (Qflex uses $M=I$, i.e.\ Euclidean steepest
ascent). Proposition~1 of the paper shows that under mild regularity the induced flow is
a \emph{valid policy-improvement flow}: the advantage
$F(t;s) = \mathbb{E}_{a\sim\pi^{(t)}}[Q] - \mathbb{E}_{a'\sim\pi^{(0)}}[Q]$
is monotone nondecreasing in $t$.

In practice the policy is a Gaussian \emph{source} $\pi^{(0)}_\theta$ composed with a
learned state-dependent velocity field $v_w$, integrated as
\begin{equation}
\pi^{(1)}_{\theta,w}(a|s,a^{(0)}) = a^{(0)} + \int_0^1 v_w(t,s,a^{(t)})\,dt .
\end{equation}
Targets $a^{(1)}$ are produced by $N=20$ finite gradient-ascent steps on the
differentiable critic, with a step size truncated by the $\ell_2$ diameter of the action
box to keep updates inside $[-1,1]^{|\mathcal{A}|}$:
\begin{equation}
a^{(\frac{n}{N})} \leftarrow a^{(\frac{n-1}{N})} + \bar\eta\,\nabla_a Q_\phi,
\qquad
\bar\eta = \min\!\left(\eta,\ \frac{2\sqrt{|\mathcal{A}|}}{\|\nabla_a Q_\phi\|}\right).
\end{equation}
The velocity field is then fit by conditional flow matching against the optimal-transport
path, $v^{(t)} = a^{(1)} - a^{(0)}$, minimising
$\mathcal{L}_v(w) = \mathbb{E}\|v_w(t,s,a^{(t)}) - (a^{(1)}-a^{(0)})\|^2$.

Two further departures from the PPO pair:
\begin{itemize}[leftmargin=1.4em,itemsep=0.2em]
\item \textbf{Off-policy.} Transitions go to a replay buffer of $10^6$ and are reused,
giving roughly two orders of magnitude better sample efficiency
($5\times10^6$ steps versus $10^9$).
\item \textbf{CrossQ-style critic.} Batch normalisation inside the $Q$-network removes the
need for a target network and permits a low update-to-data ratio. Critic quality matters
more here than in PPO, because the critic's \emph{gradient} now steers exploration, not
just the advantage estimate.
\end{itemize}

\paragraph{Learnable source distribution.} Unlike diffusion-based online RL that starts
from a fixed $\mathcal{N}(0,I)$, Qflex's source $\pi^{(0)}_\theta$ is itself trained,
which the paper reports as the main reason the flow is easy to learn. The diagnostic is
the \emph{flow superiority ratio}, the fraction of a minibatch for which
$Q(s,\pi^{(1)}) > Q(s,\pi^{(0)})$; it is reported as markedly higher on musculoskeletal
tasks than on torque-controlled ones, which is the paper's central empirical claim about
over-actuation.

\begin{table}[h!]
\centering
\small
\begin{tabular}{@{}lL{3.0cm}L{3.0cm}L{3.8cm}@{}}
\toprule
 & \textcolor{mmC}{MuscleMimic} & \textcolor{mjC}{mjlab} & \textcolor{qfC}{Qflex} \\
\midrule
Family        & on-policy PG & on-policy PG & off-policy actor--critic \\
Objective     & clipped surrogate & clipped surrogate & $Q$-maximisation + flow matching \\
Critic use    & advantage baseline & advantage baseline & \textbf{gradient steers exploration} \\
Exploration   & $\mathcal{N}(\mu,\sigma)$ & $\mathcal{N}(\mu,\sigma)$ & $Q$-guided ODE transport \\
Target net    & \na & \na & none (BatchNorm critic) \\
Sample reuse  & $E{=}1$ & $E{\approx}5$ & replay, $10^6$ \\
Baselines shipped & \na & \na & SAC, CrossQ, DACER, QSM, DIPO, QVPO, SDAC, DPMD, IDEM, DSACT \\
\bottomrule
\end{tabular}
\caption{Learning-rule comparison.}
\end{table}

% =====================================================================
\section{Task specification: what the inputs actually are}

This is where the user-facing difference is largest, and it is \emph{not} the case that
all three are ``trajectories to be followed''.

\subsection{MuscleMimic: a corpus, and a goal-conditioned generalist}

Training consumes a \emph{dataset} of retargeted clips --- 972 train / 108 test for
MyoFullBody (\texttt{KIT\_KINESIS\_TRAINING\_MOTIONS}, curated by Kinesis from AMASS+KIT),
and 1770 / 312 for MyoBimanualArm. The policy is goal-conditioned: its observation
includes the reference at $n$ future steps, so one network covers the whole corpus.

\begin{itemize}[leftmargin=1.4em,itemsep=0.2em]
\item \textbf{Goal}: \texttt{GoalTrajMimic} with \texttt{n\_step\_lookahead=5},
\texttt{n\_step\_stride=20} (0.2\,s at 100\,Hz), a normalised motion-phase scalar, and a
``concise'' encoding (root position/velocity deltas plus site relative positions) without
which 5-step lookahead over 17 sites would inflate the observation fivefold.
\item \textbf{Reward}: DeepMimic-style, $r_t = \max\{0,\, r^{\text{imit}}_t + P_t\}$ with
$r^{\text{imit}}_t = w_q r_q + w_{\dot q} r_{\dot q} + w_p r_p + w_\theta r_\theta
+ w_v r_{v\omega} + w_v r_{vv}$, each term an exponential kernel on a tracking error,
plus penalties including muscle activation energy.
\item \textbf{Termination}: mean mimic-site deviation \emph{relative to the root frame}
($\delta_{\text{site}}$), plus a world-frame pelvis check ($\delta_{\text{root}}$).
Relative rather than absolute, so global drift does not kill a biomechanically valid gait.
\item \textbf{Sampling}: clips are prioritised by EMA-smoothed early-termination rate,
with a uniform floor --- a soft replacement for Kinesis's discrete hard-negative mining.
\end{itemize}

\subsection{mjlab: commands, or one registry motion}

mjlab's two flagship tasks are structurally different from each other:
\emph{velocity tracking} is command-driven (the policy receives a sampled target velocity,
no reference motion exists), while \emph{motion tracking} follows a single preprocessed
motion fetched from a W\&B registry. Neither is a corpus-scale imitation setup. Rewards
are assembled from composable manager terms rather than a fixed formula.

\subsection{Qflex: task reward, with one reference trajectory where needed}

Most Qflex benchmarks are pure reward maximisation with no reference at all ---
\texttt{myoHandPenTwirlRandom-v0}, \texttt{h1-run-v0}, \texttt{OstrichRun-v1},
\texttt{SMPLHumanoidJump-v1}, \texttt{myoLegWalk-v0}. The MS-Human-700 locomotion task is
a hybrid, and the closest point of contact with MuscleMimic. From
\texttt{envs/ms\_700\_walk/locomotion\_v1.py}:
\begin{equation}
r = w_{\text{vel}} r_{\text{vel}} + w_{\text{qpos}} r_{\text{qpos}}
  + w_{\text{qvel}} r_{\text{qvel}} + w_{\text{act}} r_{\text{act}}
  + w_{\text{healthy}} r_{\text{healthy}}
\end{equation}
with $r_{\text{qpos}} = -\sum (q^{*} - q)^2$ over a masked joint set,
$r_{\text{qvel}} = -\|\dot q^{*} - \dot q\|$,
$r_{\text{act}} = -\|\mathbf{a}\|/n_a$ (the same muscle-energy idea as MuscleMimic and
Kinesis), and $r_{\text{vel}}$ penalising deviation from a target COM velocity
(default $v_x = 1.25$\,m/s). Reference-state initialisation is used
(\texttt{init\_qpos = qpos\_target}).

The critical structural difference: the reference is a \textbf{single} trajectory,
queried \emph{by time} through a SciPy spline
(\texttt{trajectory.query(time=self.data.time + self.\_reset\_time)}), loaded from a
\texttt{.mot} file and optionally paired with EMG. There is no clip dataset, no corpus
sampling, and no $n$-step lookahead --- only the instantaneous error
$\texttt{qpos\_dist} = q^{*} - q$ appears in the observation, together with a periodic
phase variable.

\begin{table}[h!]
\centering
\small
\begin{tabular}{@{}lL{3.0cm}L{3.0cm}L{3.6cm}@{}}
\toprule
 & \textcolor{mmC}{MuscleMimic} & \textcolor{mjC}{mjlab} & \textcolor{qfC}{Qflex (MS700)} \\
\midrule
Reference data   & 972+ clips, retargeted & 1 motion, or none & 1 \texttt{.mot} trajectory \\
Indexed by       & clip id + step & time & \textbf{wall-clock time (spline)} \\
Policy sees      & goal + 5-step lookahead & goal (tracking task) & instantaneous error + phase \\
Generalist?      & \yes, one net for corpus & per-task & per-task \\
Retargeting      & SMPL fit + GMR-Fit IK & external preprocessing & pre-baked \texttt{.mot} \\
Energy penalty   & \yes & optional term & \yes \\
RSI              & \yes & \yes & \yes \\
\bottomrule
\end{tabular}
\caption{Task-specification comparison. ``Indexed by wall-clock time'' is the sharpest
difference: Qflex's reference is a continuous function of $t$, MuscleMimic's is a
discrete corpus with adaptive sampling.}
\end{table}

% =====================================================================
\section{Technical implementation}

\subsection{Physics backends}

All three use MuJoCo, but reach it very differently.

\begin{table}[h!]
\centering
\footnotesize
\begin{tabular}{@{}lL{3.1cm}L{3.1cm}L{3.5cm}@{}}
\toprule
 & \textcolor{mmC}{MuscleMimic} & \textcolor{mjC}{mjlab} & \textcolor{qfC}{Qflex} \\
\midrule
Engine & \texttt{mujoco-mjx} 3.4.0 & \texttt{mujoco-warp} $\ge$3.8.0.3 & \texttt{gymnasium[mujoco]} 1.0.0 \\
Path to kernels & MJX with \texttt{impl='warp'} & \texttt{mujoco\_warp} directly & \textbf{CPU MuJoCo} \\
Kernel language & NVIDIA Warp (CUDA) & NVIDIA Warp (CUDA) & C \\
Also supports & MJX/XLA; CPU for eval & CPU for eval (macOS) & CPU only \\
Vectorisation & \texttt{jax.vmap}, 8192 envs & batched Warp, 4096 envs & \textbf{multiprocessing + SharedMemory, 8 envs} \\
GPU required & for training & for training & \no \\
\bottomrule
\end{tabular}
\caption{Physics. Note the first two converge on the \emph{same} Warp kernels by different
routes; Qflex does not use GPU physics at all.}
\end{table}

\noindent
The MuscleMimic/mjlab convergence deserves emphasis: because MuJoCo 3.4 unified the MJX
interface with a selectable implementation, \texttt{mjx.put\_model(model, impl='warp')}
dispatches to the same Warp kernels that \texttt{mujoco\_warp} exposes. The physics is
therefore not a fidelity difference between those two --- it is a host-language
difference. Qflex, by contrast, runs stock CPU MuJoCo behind a Gymnasium interface and
parallelises with worker processes over shared memory
(\texttt{relax/env/vector/process1.py}).

This also explains Qflex's step budget. At 8 CPU environments it cannot brute-force
$10^9$ steps, so sample efficiency is not a stylistic preference but a requirement ---
which is consistent with choosing an off-policy method.

\subsection{Learning stack}

\begin{table}[h!]
\centering
\footnotesize
\begin{tabular}{@{}lL{3.1cm}L{3.1cm}L{3.4cm}@{}}
\toprule
 & \textcolor{mmC}{MuscleMimic} & \textcolor{mjC}{mjlab} & \textcolor{qfC}{Qflex} \\
\midrule
Array framework & JAX 0.7.2 & \textbf{PyTorch $\ge$2.7} & JAX 0.4.27 \\
NN library & Flax 0.12.0 & \texttt{torch.nn} & \textbf{dm-haiku} \\
Optimiser lib & optax & torch.optim (via rsl\_rl) & optax \\
Distributions & distrax 0.1.7 & torch.distributions & hand-rolled \\
RL code & vendored ($\sim$5.5k LOC) & \texttt{rsl-rl-lib} 5.2.0 & vendored \texttt{relax} ($\sim$9.6k LOC) \\
Checkpointing & orbax-checkpoint 0.11.37 & \texttt{torch.save} & \textbf{dill} 0.3.8 \\
Config & Hydra + YAML & \texttt{tyro} dataclasses & argparse \\
Multi-GPU & \na & \texttt{torchrunx} & \na \\
Logging & W\&B & W\&B + TensorBoard & W\&B + TensorBoardX \\
\bottomrule
\end{tabular}
\caption{Learning stack. Qflex and MuscleMimic are both JAX but share no libraries:
Haiku versus Flax, and a four-minor-version JAX gap.}
\end{table}

\paragraph{Dependency versus vendoring.}
mjlab \emph{depends} on \texttt{rsl\_rl}, so it inherits every feature that library gains.
MuscleMimic and Qflex both \emph{vendor} their RL code, so they inherit nothing and
maintain everything. This single architectural choice explains the distillation gap in
Section~\ref{sec:distill} far better than JAX-versus-PyTorch does.

\subsection{Policy architectures}

\begin{table}[h!]
\centering
\footnotesize
\begin{tabular}{@{}lL{3.1cm}L{3.1cm}L{3.4cm}@{}}
\toprule
 & \textcolor{mmC}{MuscleMimic} & \textcolor{mjC}{mjlab} & \textcolor{qfC}{Qflex} \\
\midrule
Actor & $16\times1024$ gated residual, SiLU, LayerNorm & MLP (task-configured) & MLP $3\times256$ \\
Residual gate & $\sigma(g_l)$, $g_l$ init $-2$ & \na & \na \\
Extra heads & \na & \na & velocity field $3\times256$ \\
Obs normalisation & Welford running mean/std & running normaliser (rsl\_rl) & --- \\
Optimiser & \textbf{Muon} (2D) + Adam (1D) & Adam & Adam, $3\text{e-}4 \to 3\text{e-}5$ \\
Mixture of experts & Soft-MoE layers (optional) & \na & \na \\
\bottomrule
\end{tabular}
\caption{Network design. MuscleMimic's actor is by far the largest; Qflex compensates
for a small network with a better exploration rule.}
\end{table}

\noindent
The depth gap is striking --- $16\times1024$ against $3\times256$ --- and reflects the
paradigms. MuscleMimic must encode a whole motion corpus in the weights of a generalist
goal-conditioned policy; Qflex learns one task at a time and spends its capacity on the
critic and the flow instead.

\subsection{Observation pipelines}
\label{sec:obs}

\begin{itemize}[leftmargin=1.4em,itemsep=0.3em]
\item \MM. Observations are registered plug-ins with group tags. Muscle-specific types
were added to the vendored LocoMuJoCo core: \texttt{ActuatorLength},
\texttt{ActuatorVelocity}, \texttt{ActuatorForce}, \texttt{ActuatorExcitation},
\texttt{ActuatorActivation}, \texttt{TouchSensor}. Asymmetric actor/critic observation is
supported via index masking (\texttt{actor\_obs\_group} / \texttt{critic\_obs\_group}) but
is unused in the shipped configs. History stacking via
\texttt{NStepWrapper(split\_goal=True)}. \emph{No noise or latency model.}

\item \MJ. A per-term pipeline:
\texttt{compute $\rightarrow$ noise $\rightarrow$ clip $\rightarrow$ scale $\rightarrow$ delay $\rightarrow$ history}.
Noise may be stateless (\texttt{NoiseCfg}) or stateful (\texttt{NoiseModelCfg}). The delay
model is genuinely detailed: \texttt{delay\_min\_lag}/\texttt{delay\_max\_lag} sampled
per-environment, \texttt{delay\_hold\_prob} for temporally correlated latency,
\texttt{delay\_update\_period} for multi-rate sensors, \texttt{delay\_per\_env\_phase} to
destagger resampling. Named groups with \texttt{enable\_corruption} let one term
definition serve both a noisy actor group and a clean critic group.

\item \QF. Observations are assembled by the environment wrapper as a plain concatenation.
For MS700: \texttt{qpos, qvel, qacc, act, actuator\_forces} (scaled $/1000$, clipped
$\pm100$), \texttt{actuator\_length, actuator\_velocity} (clipped), \texttt{sim\_time},
a phase variable, pelvis and sternum positions, and \texttt{qpos\_dist}. Note this
includes the same muscle-state channels MuscleMimic added --- arrived at independently.
\emph{No noise, clipping, or latency framework.}
\end{itemize}

\noindent
For a controller that must eventually run on real hardware with real sensors, mjlab's
pipeline is the only one of the three that models the sensing problem.

% =====================================================================
\section{Outputs, artifacts, and deployment}
\label{sec:outputs}

\begin{table}[h!]
\centering
\footnotesize
\begin{tabular}{@{}lL{3.1cm}L{3.3cm}L{3.3cm}@{}}
\toprule
 & \textcolor{mmC}{MuscleMimic} & \textcolor{mjC}{mjlab} & \textcolor{qfC}{Qflex} \\
\midrule
Training checkpoint & Orbax OCDBT \emph{directory} & \texttt{.pt} (\texttt{torch.save}) & \texttt{.pkl} (dill) \\
Contents & weights, optimiser, Welford stats, RNG, \texttt{\_sharding} & state dicts + infos & params + Haiku state \\
Deployment artifact & \textbf{none} & \textbf{\texttt{.onnx}, opset 18} & \texttt{stochastic.pkl} / \texttt{deterministic.pkl} \\
Normaliser baked in & \na & \yes & \no \\
Embedded metadata & \na & \yes (see below) & \texttt{.txt} sidecar \\
Runs without the training framework & \no & \yes & \no \\
ONNX packages installed & \textbf{0} & 8 & 0 \\
\bottomrule
\end{tabular}
\caption{Output formats. Verified: \texttt{find . -name '*.onnx'} returns zero results in
both MuscleMimic and Qflex, and neither declares an ONNX dependency.}
\end{table}

\subsection{MuscleMimic: a training checkpoint, not a policy}

\begin{lstlisting}
checkpoints/<run-hash>/checkpoint_8500/
  _CHECKPOINT_METADATA
  train_state/            # OCDBT / tensorstore
    _sharding             # JAX device-sharding spec
    manifest.ocdbt
    array_metadatas/
    ocdbt.process_0/
  config/metadata         # JSON: full Hydra config
  metadata/metadata       # JSON: step, update_number, lr, ...
\end{lstlisting}

\noindent
Running a policy requires JAX, the \texttt{musclemimic} package, and reconstruction of the
MuJoCo environment object to recover observation indexing
(\texttt{play\_policy} / \texttt{play\_policy\_mujoco}). There is no export step of any
kind --- no ONNX, TorchScript, SavedModel, StableHLO, or TFLite.

\subsection{mjlab: a self-describing deployment artifact}

ONNX export runs automatically on every checkpoint save. The exported graph bakes in the
observation normaliser and the deterministic output head, so the artifact is
raw-observations $\rightarrow$ actions with no external state. Metadata is written into
\texttt{model.metadata\_props}:

\begin{lstlisting}
run_path, joint_names, joint_stiffness, joint_damping,
default_joint_pos, command_names, observation_names, action_scale
\end{lstlisting}

\noindent
The motion-tracking variant goes further, registering the reference motion itself as ONNX
buffers (\texttt{joint\_pos}, \texttt{joint\_vel}, \texttt{body\_pos\_w}, \ldots) and
exposing a \texttt{(obs, time\_step)} signature, so a tracking policy ships as one file
with no Python-side trajectory loader. TorchScript export is also available via
\texttt{\_TorchMLPModel}.

\paragraph{One caveat for muscle models.} \texttt{get\_base\_metadata()} asserts
\texttt{isinstance(joint\_action, JointPositionAction)} and reads
\texttt{actuator\_gainprm}/\texttt{actuator\_biasprm} as PD gains. A muscle-actuated model
has neither, so this helper must be given a muscle variant before the ONNX pipeline works
for an MSK embodiment.

\subsection{Qflex: pickled JAX callables}

\texttt{save\_policy\_structure} uses a custom dill-based pickler to persist \emph{traced
callables} (\texttt{stochastic.pkl}, \texttt{deterministic.pkl}) alongside human-readable
\texttt{.txt} summaries; \texttt{save\_policy} writes \texttt{(params, state)} tuples
periodically, and \texttt{state.pkl} holds the full algorithm state at the end. This is
portable across machines but not across frameworks, and carries the usual caveat that
unpickling executes code.

\subsection{External tool compatibility}

\begin{table}[h!]
\centering
\small
\begin{tabular}{@{}lccc@{}}
\toprule
Consumer & \textcolor{mmC}{MuscleMimic} & \textcolor{mjC}{mjlab} & \textcolor{qfC}{Qflex} \\
\midrule
ONNX Runtime (C/C++, ARM)      & \no & \yes & \no \\
TensorRT                        & \no & \yes & \no \\
TorchScript / LibTorch          & \no & \yes & \no \\
Netron graph inspection         & \no & \yes & \no \\
ROS\,2 / real-time C++ node     & \no & \yes & \no \\
Reload into own framework       & \yes & \yes & \yes \\
HuggingFace hosting             & \yes & \yes & \na \\
\bottomrule
\end{tabular}
\caption{Only mjlab produces an artifact consumable outside its own Python process.}
\end{table}

\noindent
The practical obstacle is not the weights. All three policies are MLPs, and extracting
arrays into NumPy is straightforward. The obstacle is the \textbf{observation contract}:
which quantities, in which order, with what history stacking, goal lookahead, and
normalisation. mjlab records this inside the artifact; the other two require the training
framework to reconstruct it.

% =====================================================================
\section{Musculoskeletal model support}
\label{sec:msk}

\begin{table}[h!]
\centering
\footnotesize
\begin{tabular}{@{}lL{3.0cm}L{3.0cm}L{3.5cm}@{}}
\toprule
 & \textcolor{mmC}{MuscleMimic} & \textcolor{mjC}{mjlab} & \textcolor{qfC}{Qflex} \\
\midrule
MSK assets shipped & MyoFullBody, MyoBimanualArm & \textbf{none} (G1, Go1, YAM) & MS-Human-700 (submodule), MyoSuite \\
Largest model & 416 muscles, 123 joints, 72 DoF & \na & \textbf{700 muscles, 85 DoF} \\
Muscle actuator support & native & \texttt{BuiltinMuscleActuator}, \texttt{XmlActuatorCfg} & via MyoSuite/MuJoCo \\
Tendon transmission & \yes & \yes (\texttt{TendonEffortAction}) & \yes \\
Muscle-state obs & 5 channels + touch & user-defined terms & 3 channels \\
Activation-energy penalty & \yes & user-defined & \yes \\
Validated against & kinematics, kinetics, GRF, EMG & \na & task reward only \\
\bottomrule
\end{tabular}
\caption{Muscle support. Measured directly: \texttt{MS-Human-700.xml} loads as
\texttt{nq=85, nu=na=700, ntendon=700}; the bundled
\texttt{MS-Human-700-Locomotion.xml} subset is \texttt{nq=36, nu=100}, and the
environment loads the full model.}
\end{table}

\noindent
mjlab has no biomechanical assets but is architecturally ready for them. Its
\texttt{XmlActuatorCfg} docstring states explicitly that MJCF containing
\texttt{<muscle>} elements is wrapped without re-specifying gains, and
\texttt{mujoco\_warp} implements the full Hill-type stack in Warp kernels
(\texttt{muscle\_dynamics}, \texttt{muscle\_gain}, \texttt{muscle\_bias}, and the
corresponding \texttt{DynType}/\texttt{GainType}/\texttt{BiasType} enums). Importing an
MSK model is therefore a matter of porting the MDP, not the physics.

\paragraph{No first-party musculoskeletal result.} To be precise about what mjlab has and
has not shown: its twelve registered tasks are \texttt{Mjlab-Cartpole-\{Balance,Swingup\}},
four YAM manipulation tasks, two Unitree G1 tracking tasks, and four velocity-tracking
tasks on G1/Go1. None involves a muscle-driven model, and \texttt{asset\_zoo} contains only
\texttt{unitree\_g1}, \texttt{unitree\_go1}, and \texttt{i2rt\_yam}. The sole musculoskeletal
connection is a \emph{third-party} entry on mjlab's ``Projects built on mjlab'' page,
\texttt{MyoHub/mjlab\_myosuite} (``Musculoskeletal simulation integration with MyoSuite''),
which is neither vendored nor a dependency --- \texttt{grep} for it across \texttt{src/},
\texttt{pyproject.toml}, and \texttt{uv.lock} returns nothing. It is assessed separately in
Section~\ref{sec:mjlabmyo}.

% =====================================================================
\section{MyoSuite: the common substrate}
\label{sec:myosuite}

MyoSuite is not a fourth competitor; it is the layer the other three stand on. MuscleMimic's
embodiments are assembled from MyoSuite components (MyoArm, MyoLegs, MyoBack); Qflex
benchmarks on \texttt{myoLegWalk-v0} and \texttt{myoHandPenTwirlRandom-v0}; and
\texttt{mjlab\_myosuite} (Section~\ref{sec:mjlabmyo}) exists solely to bridge to it.
Understanding what a MyoSuite environment \emph{is} therefore clarifies several earlier
sections. Measurements below are against \textbf{v2.12.2}, which pins
\texttt{mujoco>=3.6,<3.7}.

\subsection{Anatomy of an environment}

A MyoSuite environment is one registration entry, one MuJoCo XML, and one Python class.

\begin{lstlisting}
register(
    id="myoLegWalk-v0",
    entry_point="myosuite.envs.myo.myobase.walk_v0:WalkEnvV0",
    max_episode_steps=1000,
    kwargs={"model_path": curr_dir + leg_model, "normalize_act": True,
            "min_height": 0.8, "max_rot": 0.8, "hip_period": 100,
            "reset_type": "init", "target_x_vel": 0.0, "target_y_vel": 1.2,
            "target_rot": None},
)
\end{lstlisting}

\noindent
Task variants are simply different \texttt{kwargs}. Models live in
\texttt{myosuite/simhive/}, a hive of asset submodules (\texttt{myo\_sim},
\texttt{MPL\_sim}, \texttt{object\_sim}, \texttt{YCB\_sim},
\texttt{furniture\_sim}). The class sits three levels deep:

\begin{lstlisting}
gym.Env + EzPickle + ObsVecDict
  \-- MujocoEnv      (envs/env_base.py)      generic MuJoCo + the MDP contract
       \-- BaseV0    (envs/myo/base_v0.py)   muscle-specific behaviour
            \-- WalkEnvV0 (myobase/walk_v0.py)  the task
\end{lstlisting}

\subsection{The MDP contract: two lists and two methods}

\texttt{MujocoEnv.\_setup} fixes the contract --- \texttt{obs\_keys} (which entries of an
observation dict form the policy input), \texttt{weighted\_reward\_keys} (keys and weights
that sum to the reward), plus \texttt{proprio\_keys}, \texttt{visual\_keys},
\texttt{reward\_mode} (\texttt{dense}/\texttt{sparse}), \texttt{frame\_skip},
\texttt{normalize\_act}, and \texttt{obs\_range} (default $\pm 10$). A concrete task then
supplies two lists and two methods:

\begin{lstlisting}
DEFAULT_OBS_KEYS = ["qpos_without_xy", "qvel", "com_vel", "torso_angle",
    "feet_heights", "height", "feet_rel_positions", "phase_var",
    "muscle_length", "muscle_velocity", "muscle_force"]

DEFAULT_RWD_KEYS_AND_WEIGHTS = {"vel_reward": 5.0, "done": -100,
    "cyclic_hip": -10, "ref_rot": 10.0, "joint_angle_rew": 5.0}

def get_obs_dict(self, mj_model, mj_data) -> dict         # named quantities
def get_reward_dict(self, obs_dict) -> OrderedDict        # named components
\end{lstlisting}

\noindent
The observation vector is \texttt{concatenate(obs\_dict[k] for k in obs\_keys)}.
\texttt{get\_reward\_dict} must emit three keys --- \texttt{"sparse"}, \texttt{"solved"},
and \texttt{"done"} --- after which \texttt{rwd\_dict["dense"]} is computed as
$\sum_k w_k\,\mathrm{rwd}_k$. Three consequences follow, and all three were already visible
in earlier sections:

\begin{enumerate}[leftmargin=1.4em,itemsep=0.25em]
\item \textbf{Muscle state is already first-class.} \texttt{muscle\_length},
\texttt{muscle\_velocity} and \texttt{muscle\_force} are default observations, and
\texttt{BaseV0} auto-appends \texttt{"act"} whenever \texttt{mj\_model.na > 0}. MuscleMimic's
\texttt{ActuatorLength}/\texttt{ActuatorVelocity}/\texttt{ActuatorForce} observation types
(Section~\ref{sec:msk}) are the same channels, arrived at independently.
\item \textbf{There is no processing layer.} The path from \texttt{get\_obs\_dict} to the
network is concatenation and a clip to \texttt{obs\_range}. No noise, scaling, delay, or
history --- no seam at which to attach them. This is the structural reason wrapping MyoSuite
cannot recover mjlab's observation pipeline (Section~\ref{sec:obs}).
\item \textbf{Termination lives inside the reward.} \texttt{"done"} is a key of
\texttt{rwd\_dict}, so reward and episode end are computed in one method rather than by
independent managers.
\end{enumerate}

\paragraph{Clinical conditions.} \texttt{BaseV0} exposes a \texttt{muscle\_condition} switch
with no analogue in any of the three frameworks compared here:

\begin{lstlisting}
"sarcopenia"      -> halve every actuator's peak force (gainprm[:,2] *= 0.5)
"fatigue"         -> attach the 3CC-r cumulative fatigue model
"reafferentation" -> simulate EIP -> EPL tendon-transfer surgery
\end{lstlisting}

\noindent
plus \texttt{fatigue\_reset\_vec} and \texttt{fatigue\_reset\_random}. For prosthetics and
rehabilitation work this is a substantive capability, not a curiosity.

\subsection{Training: adapters, not algorithms}

MyoSuite implements \textbf{no RL algorithm of its own}. Its runtime dependencies contain no
RL library at all (\texttt{gymnasium}, \texttt{mujoco}, \texttt{h5py}, \texttt{numpy},
\texttt{pink-noise-rl}, \ldots). What it ships instead is integration:

\begin{table}[h!]
\centering
\small
\begin{tabular}{@{}L{4.4cm}L{8.4cm}@{}}
\toprule
Component & What it wraps \\
\midrule
\texttt{mjrl\_job\_script.py}, \texttt{hydra\_mjrl\_launcher.py} & \textbf{MJRL} --- Natural Policy Gradient; the headline baseline \\
\texttt{sb3\_job\_script.py}, \texttt{hydra\_sb3\_launcher.py} & \textbf{Stable-Baselines3} --- PPO and SAC configs \\
\texttt{torchrl\_job\_script.py} & \textbf{TorchRL} \\
\texttt{baseline\_Reflex/reflexCtr.py} & a non-learned reflex controller \\
\texttt{SAR\_pretrained/} & pretrained Synergistic Action Representation policies \\
tutorials \texttt{4a}, \texttt{4c}, \texttt{SAR/} & \textbf{DEP-RL}, SB3, SAR walkthroughs \\
\bottomrule
\end{tabular}
\caption{\texttt{myosuite/agents/} is a baselines directory: only 8 Python files
($\approx$2{,}000 lines) among 361, the remainder being YAML configs, pickled policies,
CSV logs and plots. Hydra and submitit provide cluster launching.}
\end{table}

\subsection{GPU training: a second, disjoint implementation}

MyoSuite does have a GPU path, in main under the \texttt{[mjx]} extra --- but it is not an
acceleration of the classic environments. It is a \emph{parallel class hierarchy} that
shares no base class, no method names, and no array library with them:

\begin{table}[h!]
\centering
\footnotesize
\begin{tabular}{@{}L{3.0cm}L{4.6cm}L{4.8cm}@{}}
\toprule
 & classic path & MJX path \\
\midrule
Base class   & \texttt{gym.Env} $\to$ \texttt{MujocoEnv} $\to$ \texttt{BaseV0} & \texttt{mjx\_env.MjxEnv} (MuJoCo Playground) $\to$ \texttt{MjxMyoBase} \\
Gym import   & \yes & \textbf{none} \\
Reset        & \texttt{reset(**kwargs)} & \texttt{reset(rng) -> State} \\
Step         & \texttt{step(a)}, mutates \texttt{self.mj\_data} & \texttt{step(state, action) -> State}, \textbf{pure} \\
Task hooks   & \texttt{get\_obs\_dict}, \texttt{get\_reward\_dict} & \texttt{\_get\_obs}, \texttt{\_get\_rewards}, \texttt{\_get\_done}, \texttt{\_get\_metrics} \\
Arrays       & NumPy & \texttt{jax.numpy} \\
Trainer      & MJRL / SB3 / TorchRL / DEP-RL & \textbf{Brax PPO} \\
\textbf{Environments} & \textbf{31} & \textbf{2} (\texttt{playground\_pose\_v0}, \texttt{playground\_reach\_v0}) \\
\bottomrule
\end{tabular}
\caption{The two MyoSuite paths. \texttt{MjxPoseEnvV0} and \texttt{PoseEnvV0} are the same
task written twice.}
\end{table}

\noindent
Every line of \texttt{MjxMyoBase.step} is a \texttt{state.replace(...)} --- a new pytree,
never a mutation of \texttt{self}. That purity is the whole point, because MyoSuite performs
no batching itself; it hands the env to Brax:

\begin{lstlisting}
from mujoco_playground import wrapper
make_inference_fn, params, _ = ppo.train(
    num_envs=env._config.num_envs,
    wrap_env_fn=wrapper.wrap_for_brax_training, ...)
\end{lstlisting}

\noindent
Playground's wrapper batches, and Brax's PPO \texttt{jit}s and \texttt{vmap}s across
environments. The shipped \texttt{ppo\_config} is 50M timesteps, lr $3\times10^{-4}$,
$\gamma=0.97$, $\lambda=0.95$, $\epsilon=0.3$, 32 minibatches, \textbf{8} update epochs,
batch 256, unroll 10, normalised observations, and $(64,64,64)$ policy and value networks
--- note that this is a far smaller network and a far more aggressive epoch count than
MuscleMimic's $16\times1024$ at $E=1$ (Section~\ref{sec:distill}).

\paragraph{Why the coverage gap is structural.} A Gymnasium environment is a stateful Python
object owning one \texttt{mj\_data}; \texttt{jax.vmap} requires a pure function over array
pytrees. No wrapper converts one into the other --- the only options are running $N$ copies
in a loop or rewriting the environment. Hence 31 classic environments and 2 MJX ones: the
rewrite is manual and per-environment. This single fact explains both MyoSuite's $\approx$6\%
GPU coverage and the vectorisation behaviour of \texttt{mjlab\_myosuite} below.

\begin{table}[h!]
\centering
\small
\begin{tabular}{@{}L{4.2cm}L{4.4cm}L{4.0cm}@{}}
\toprule
 & Environment model & Batching performed by \\
\midrule
MyoSuite (classic)   & stateful Gym object & nothing built in \\
MyoSuite (MJX)       & pure JAX function & Brax / Playground \\
\textcolor{mmC}{MuscleMimic} & pure JAX function & its own \texttt{vmap} \\
\textcolor{mjC}{mjlab}       & batched Warp + torch tensors & \textbf{the env itself} \\
\textcolor{qfC}{Qflex}       & stateful Gym object & 8 CPU processes (shared memory) \\
\bottomrule
\end{tabular}
\caption{The fault line running through every project here. MyoSuite-MJX and MuscleMimic are
architecturally the same kind of object; mjlab takes the third road, batching inside the
environment rather than replicating a pure function.}
\end{table}

% =====================================================================
\section{Addendum: \texttt{mjlab\_myosuite}, a third-party MSK bridge}
\label{sec:mjlabmyo}

Since mjlab ships no biomechanical assets, the obvious question is whether the community
project referenced in Section~\ref{sec:msk} closes that gap. It does, partially, and in a
way that is worth understanding precisely --- because the architecture it chose forfeits
most of the mjlab features that motivated the comparison in the first place.

\subsection{What it is}

\texttt{MyoHub/mjlab\_myosuite} is a \emph{fork} of mjlab by git history --- 280 commits
reaching back to mjlab's ``Initial public release'' of 2025-09-29 --- that was later
restructured to consume mjlab as an ordinary dependency (commit \texttt{9d31812},
``use mjlab as dependency''). Its own source tree is 4{,}899 lines across 19 modules, plus
1{,}615 lines of tests. The README opens with a banner reading
\emph{``Early Prototype --- breaking changes are likely''}, and the package metadata
declares \texttt{Development Status :: 4 - Beta} at version \texttt{0.1.0}.

\subsection{Architecture: a wrapper, not a port}

This is the decisive design fact. \texttt{mjlab\_myosuite} does \textbf{not} bring
musculoskeletal models into mjlab's manager-based MDP. It leaves the MDP inside MyoSuite
and fabricates just enough of mjlab's interface for mjlab's training scripts to run
against it. The bulk of the package --- \texttt{wrapper.py}, 1{,}457 lines --- defines:

\begin{lstlisting}
class _MockActionManager: ...
class _MockObservationManager: ...
class _MockCommandManager: ...
class _MockScene: ...
class MyoSuiteVecEnvWrapper(VecEnv, gym.Env): ...
\end{lstlisting}

\noindent
So a MyoSuite Gym environment is adapted to \texttt{rsl\_rl}'s \texttt{VecEnv} and dressed
in mock managers that satisfy attribute probes such as
\texttt{hasattr(env, "action\_manager")}.

\begin{table}[h!]
\centering
\small
\begin{tabular}{@{}L{5.6cm}cc@{}}
\toprule
mjlab capability & native mjlab task & via \texttt{mjlab\_myosuite} \\
\midrule
PPO training loop (\texttt{rsl\_rl})        & \yes & \yes \\
W\&B logging, checkpointing                 & \yes & \yes \\
ONNX export                                 & \yes & partial (see below) \\
Viser / native viewer playback              & \yes & \yes \\
\midrule
Observation noise, clipping, scaling        & \yes & \textbf{\no} \\
Sensor-latency model (\texttt{delay\_*})    & \yes & \textbf{\no} \\
Observation history stacking                & \yes & \textbf{\no} \\
Reward manager (composable terms)           & \yes & \textbf{\no} \\
Event manager / domain randomisation        & \yes & \textbf{\no} \\
Curriculum manager                          & \yes & \textbf{\no} \\
\texttt{BuiltinMuscleActuator} / \texttt{XmlActuatorCfg} & \yes & \textbf{\no} \\
Terrain, scene composition                  & \yes & \textbf{\no} \\
\bottomrule
\end{tabular}
\caption{What survives the bridge. Everything below the rule is a manager-layer feature,
and the manager layer is exactly what the mock objects replace. The MDP --- observations,
rewards, terminations, actuation --- remains MyoSuite's.}
\end{table}

\noindent
For the deployment use case this is the crux: the single strongest argument for routing
musculoskeletal work through mjlab was its per-term noise and latency pipeline
(Section~\ref{sec:obs}), and that pipeline is precisely what the wrapper bypasses.

\paragraph{It contributes no parallelism.} Given the two disjoint MyoSuite paths of
Section~\ref{sec:myosuite}, the wrapper's vectorisation behaviour follows directly
(\texttt{wrapper.py}:190--207):

\begin{lstlisting}
if isinstance(env, vector.VectorEnv):
    self.env = env                 # already batched (MyoSuite MJX) -> pass through
    self.num_envs = env.num_envs
else:
    env_factories = [lambda: copy.deepcopy(env_instance)
                     for _ in range(num_envs)]
    self.env = SyncVectorEnv(env_factories)     # <-- serial Python loop
\end{lstlisting}

\noindent
For a \emph{classic} MyoSuite environment --- the stateful kind, and the only kind that
covers locomotion --- the result is \texttt{SyncVectorEnv}: single-threaded, one environment
stepped after another, not even multiprocessing. That is plausibly slower than MyoSuite's own
subprocess-based baselines. The README's ``MJX/Warp support'' means only that an
\emph{already-batched} MyoSuite MJX environment is passed through untouched --- and there are
two such environments, neither of them locomotion. Either way, mjlab contributes zero
batching: the GPU parallelism, where it exists, belongs to Brax and Playground.

\subsection{What it adds over MyoSuite alone}

Because MyoSuite already ships a GPU PPO path, W\&B logging, checkpointing, cluster
launchers and pretrained baselines (Section~\ref{sec:myosuite}), the honest question is not
what \texttt{mjlab\_myosuite} provides but what it provides \emph{that MyoSuite did not
already have}.

\begin{table}[h!]
\centering
\small
\begin{tabular}{@{}L{5.0cm}L{3.6cm}L{3.8cm}@{}}
\toprule
Capability & MyoSuite already has & \texttt{mjlab\_myosuite} adds \\
\midrule
PPO on GPU              & \yes\ Brax + Playground & rsl-rl PPO instead \\
PPO / SAC on CPU        & \yes\ SB3, MJRL, TorchRL & \na \\
W\&B logging, checkpoints & \yes & \na \\
Hydra / cluster launchers & \yes\ submitit & \na \\
Pretrained baselines    & \yes\ SAR, NPG, DEP-RL & \na \\
\midrule
ONNX export             & \no & \textbf{new} \\
Viser playback          & \no & \textbf{new} \\
rsl-rl as the learner   & \no & \textbf{new} \\
Unified \texttt{train}/\texttt{play} CLI & shell scripts & tidier \\
\bottomrule
\end{tabular}
\caption{The delta. Three genuine additions; everything above the rule was already present.}
\end{table}

\begin{itemize}[leftmargin=1.4em,itemsep=0.25em]
\item \textbf{Automatic registration} of MyoSuite environments as
\texttt{Mjlab-MyoSuite-<env\_id>}, with a heuristic that treats an id containing
\texttt{myodm} as a tracking environment (the source comments acknowledge this is
provisional).
\item \textbf{Dual-backend detection}: accepts either a classic MyoSuite environment or an
already-batched MJX one, deciding at runtime --- though as shown above only the second is
parallel, and only two such environments exist.
\item \textbf{An ONNX metadata fallback} --- a genuine answer to the
\texttt{get\_base\_metadata} problem flagged in Section~\ref{sec:msk}.
\texttt{attach\_myosuite\_onnx\_metadata} probes whether the environment is a real
\texttt{ManagerBasedRlEnv} (checking \texttt{scene["robot"].spec.actuators}), and on
failure emits \texttt{\_get\_myosuite\_metadata} instead: action and observation
dimensions only. That avoids the assertion, but the resulting artifact carries none of
the joint names, gains, or default pose that make mjlab's native export self-describing.
\item \textbf{A tracking-task scaffold} under \texttt{tasks/tracking/}, which the README
labels ``not yet fully implemented''.
\item \textbf{Viser playback utilities} (\texttt{viewer\_helpers.py}, 486 lines).
\end{itemize}

\subsection{Version compatibility: it does not run against current mjlab}

\begin{table}[h!]
\centering
\small
\begin{tabular}{@{}lL{4.0cm}L{4.0cm}@{}}
\toprule
Dependency & pinned by \texttt{mjlab\_myosuite} & current \\
\midrule
mjlab (\texttt{pyproject}) & \texttt{git+\ldots/mjlab.git} --- \textbf{no rev, no version bound} & \na \\
mjlab (\texttt{uv.lock})   & commit \texttt{46f17fbe}, 2026-01-05 & \textbf{v1.4.0}, 2026-05-26 \\
\texttt{rsl-rl-lib}        & \textbf{3.1.0} (2025-09-18) & \textbf{5.2.0} \\
\texttt{mujoco-warp}       & rev \texttt{e605c406} & rev \texttt{88b55fc2} \\
MyoSuite                   & \texttt{git+\ldots/myosuite.git@mjx} (branch) & \na \\
Last commit                & 2026-02-05 & \na \\
\bottomrule
\end{tabular}
\caption{Dependency pins. The locked mjlab commit predates even mjlab's own v1.0.0 release
(2026-01-28) by three weeks, and the project has not been touched since 2026-02-05.}
\end{table}

\noindent
Resolving all 52 distinct \texttt{mjlab}/\texttt{rsl\_rl} symbols that
\texttt{mjlab\_myosuite} imports against an installed mjlab 1.4.0 and
\texttt{rsl-rl-lib} 5.2.0 yields the following genuine failures (import paths that merely
moved but remain re-exported are excluded, since the code has fallback chains that catch
them):

\begin{table}[h!]
\centering
\footnotesize
\begin{tabular}{@{}L{5.4cm}L{3.0cm}L{4.4cm}@{}}
\toprule
Import & Site & Status in current versions \\
\midrule
\texttt{mjlab.utils.lab\_api.rl.exporter\allowbreak.export\_policy\_as\_onnx}
  & \texttt{rl/exporter.py}
  & module tree deleted by mjlab commit \texttt{7afd6303} (``Replace custom ONNX exporter with rsl-rl 4.0 \texttt{as\_onnx()}'') \\
\texttt{mjlab.rl.exporter\_utils\allowbreak.export\_policy\_as\_onnx}
  & \texttt{rl/exporter.py} (fallback)
  & absent; that module defines only \texttt{list\_to\_csv\_str}, \texttt{get\_base\_metadata}, \texttt{attach\_metadata\_to\_onnx} \\
\texttt{mjlab.viewer.viser.conversions}
  & \texttt{viewer\_helpers.py}, \texttt{scripts/play.py}
  & module absent in 1.4.0 \\
\texttt{self.logger\_type}
  & \texttt{rl/runner.py}
  & moved to \texttt{self.logger.logger\_type} in rsl-rl 5.x \\
\texttt{self.disable\_logs}
  & \texttt{rl/runner.py}
  & moved to \texttt{self.logger.disable\_logs} \\
\texttt{self.alg.policy}
  & \texttt{rl/runner.py}
  & rsl-rl 5.x exposes \texttt{self.alg.get\_policy()} \\
\texttt{actor\_obs\_normalization}, \texttt{actor\_obs\_normalizer}
  & \texttt{rl/runner.py}
  & renamed to \texttt{obs\_normalization} / \texttt{obs\_normalizer} on \texttt{MLPModel} \\
\bottomrule
\end{tabular}
\caption{Verified incompatibilities with mjlab 1.4.0 / rsl-rl-lib 5.2.0.}
\end{table}

\paragraph{The failure is not graceful.} Both ONNX import paths fail, so
\texttt{export\_policy\_as\_onnx} degrades to a stub that raises \texttt{ImportError} ---
recoverable, since the caller catches it. More seriously,
\texttt{MyoSuiteOnPolicyRunner.save()} guards its export block with

\begin{lstlisting}
super().save(path, infos)
if self.logger_type in ["wandb"]:   # <-- outside the try/except
  try:
    ...
\end{lstlisting}

\noindent
Because \texttt{logger\_type} is read \emph{outside} the surrounding
\texttt{try}, and no longer exists on \texttt{OnPolicyRunner} in rsl-rl 5.x, the first
checkpoint save raises \texttt{AttributeError} rather than skipping export. Training would
fail at the first save interval, not at startup.

\subsection{Assessment}

\begin{itemize}[leftmargin=1.4em,itemsep=0.3em]
\item \textbf{It is a learner swap with an ONNX bolt-on, not an infrastructure upgrade.}
The implicit premise --- that MyoSuite lacks GPU training infrastructure and mjlab supplies
it --- does not hold: MyoSuite already had a Brax/Playground MJX trainer with W\&B and
published benchmarks. What actually changes is Brax PPO $\to$ rsl-rl PPO, plus viewers and
export, while parallelism is unchanged or worse.

\item \textbf{Not usable as-is} against mjlab 1.4.0. It is pinned to a pre-1.0.0 mjlab
commit and to \texttt{rsl-rl-lib} 3.1.0, two major versions behind. Because
\texttt{pyproject.toml} declares mjlab without a version bound while \texttt{uv.lock}
pins a commit, a fresh \texttt{uv sync} reproduces the old pin but \texttt{pip install -e .}
would pull current mjlab and break. Reviving it means either pinning mjlab to
\texttt{46f17fbe} (losing 5 months of mjlab development, including the muscle-actuator and
delay-model work) or porting roughly seven call sites forward.

\item \textbf{Architecturally it solves a different problem than the one posed here.} It
brings MyoSuite's environments to mjlab's \emph{trainer}. It does not bring musculoskeletal
models into mjlab's \emph{MDP}, so it confers none of the observation-corruption,
latency, randomisation, or curriculum machinery. For a project whose goal is a deployable
controller under realistic sensing, that is the wrong half of mjlab to inherit.

\item \textbf{It is nonetheless useful prior art} in two specific respects: the
\texttt{VecEnv} adaptation pattern in \texttt{wrapper.py}, and the ONNX metadata fallback,
which is the only existing demonstration of exporting a muscle-actuated policy through
mjlab's tooling. It is also the only route to \texttt{rsl\_rl} specifically --- relevant
because \texttt{rsl\_rl} is the sole stack among all those compared here that ships
student/teacher distillation (Section~\ref{sec:distill}).

\item \textbf{A native port remains the better path} for MSK work that wants mjlab's
strengths: load the MJCF through \texttt{XmlActuatorCfg} with tendon transmission, drive it
with \texttt{TendonEffortAction}, and express the imitation MDP as mjlab manager terms.
That keeps the observation pipeline and the self-describing ONNX export, at the cost of
rewriting the task definition.
\end{itemize}

% =====================================================================
\section{Scale and compute}

\begin{table}[h!]
\centering
\small
\begin{tabular}{@{}lL{2.9cm}L{2.9cm}L{3.4cm}@{}}
\toprule
 & \textcolor{mmC}{MuscleMimic} & \textcolor{mjC}{mjlab} & \textcolor{qfC}{Qflex} \\
\midrule
Parallel envs & 8192 (GPU) & 4096 (GPU) & 8 (CPU processes) \\
Throughput & $1.3\times10^4$ steps/s (H100) & comparable & CPU-bound \\
Total steps & $10^9$ & $\sim\!10^8$--$10^9$ & $5\times10^6$ \\
Wall clock & $\approx$20\,h, 1 GPU & minutes--hours & hours \\
Batch structure & 20 steps $\times$ 8192 envs, 128 minibatches & configurable & 256 from replay \\
Hardware floor & NVIDIA GPU & NVIDIA GPU & \textbf{none} \\
\bottomrule
\end{tabular}
\caption{Compute. The $200\times$ step-budget gap between MuscleMimic and Qflex is the
clearest quantitative expression of on-policy throughput versus off-policy efficiency.}
\end{table}

% =====================================================================
\section{Implications for distillation and deployment}
\label{sec:distill}

For a workflow that ends in a prosthesis-specific controller running on embedded hardware
with limited, noisy, delayed observations, the three differ sharply.

\begin{enumerate}[leftmargin=1.4em,itemsep=0.4em]
\item \textbf{Only mjlab has a distillation path.} \texttt{rsl\_rl} 5.2.0 ships
\texttt{algorithms/distillation.py} (student/teacher \texttt{MLPModel}, MSE or Huber loss,
\texttt{gradient\_length} for BPTT through recurrent students) and
\texttt{runners/distillation\_runner.py}, together with \texttt{resolve\_obs\_groups} for
giving student and teacher genuinely different observation sets. mjlab itself does
\emph{not} wire this up --- the only references are a changelog line and an aspirational
sentence in \texttt{motivation.rst} --- so it is a thin subclass plus config, not a
working feature. MuscleMimic and Qflex have nothing.

\item \textbf{Only mjlab models realistic sensing.} A prosthesis controller must tolerate
encoder quantisation, IMU noise, and variable latency. mjlab's per-term noise and delay
configuration expresses exactly this. MuscleMimic can restrict \emph{which} observations
the actor sees, but cannot corrupt or delay them.

\item \textbf{Only mjlab emits a deployable artifact,} as established in
Section~\ref{sec:outputs}.

\item \textbf{Qflex is the most interesting algorithm for the hardest regime.} If the
binding constraint becomes exploration in a 400--700 dimensional muscle action space
rather than throughput --- for instance when learning a novel prosthetic gait with no
reference motion --- Qflex's value-guided flow is the only one of the three that
addresses it directly, and it does so without dimensionality reduction. Its own results
report the advantage widening precisely as over-actuation increases.

\item \textbf{MuscleMimic remains the best teacher-producer.} It is the only one with
validated MSK embodiments, a retargeted corpus, and a generalist policy --- which is what
one distils \emph{from}.
\end{enumerate}

\paragraph{A plausible division of labour.} Train the generalist musculoskeletal teacher
in MuscleMimic, where the models, datasets, and reward shaping already work. Use mjlab as
the distillation and deployment target for prosthesis-specific and human-specific students
under realistic observation models. That reduces the JAX$\rightarrow$PyTorch boundary to a
one-time weight-and-contract transfer for an MLP, rather than a reimplementation of the
training stack. Borrow from Qflex if and when exploration, not throughput, becomes the
bottleneck.

\paragraph{Cheapest decisive experiment.} Load an existing MSK model into a bare mjlab
\texttt{Entity} with \texttt{XmlActuatorCfg(transmission\_type=TENDON)}, step it with zero
actions, and compare a few hundred steps of \texttt{qpos} against MuscleMimic's CPU MuJoCo
path. That single test exercises the MuJoCo version gap (3.4 versus 3.8), muscle-kernel
agreement, and the contact budget together.

% =====================================================================
\appendix
\section{Verification}

Every claim above was checked against the source trees rather than documentation.
Representative commands:

\begin{lstlisting}
# mjlab has no JAX
grep -rn "import jax" mjlab/src mjlab/tests        # zero hits
ls mjlab/.venv/lib/python3.13/site-packages/jax    # does not exist

# MuscleMimic and Qflex have no ONNX
find musclemimic -name '*.onnx' -not -path './.venv/*' | wc -l   # 0
find Qflex       -name '*.onnx' | wc -l                          # 0

# PureJaxRL lineage in the PPO pair
grep -rn "unroll=16"         loco-mujoco/loco_mujoco/algorithms/
grep -rn "_update_minbatch"  loco-mujoco/loco_mujoco/algorithms/

# rsl_rl distillation exists but mjlab never calls it
ls  mjlab/.venv/lib/python3.13/site-packages/rsl_rl/algorithms/distillation.py
grep -rn "Distillation" mjlab/src                  # no code hits

# Qflex physics is CPU MuJoCo
grep -rn "mjx\|mujoco_warp\|warp" Qflex/relax Qflex/requirements   # zero hits
grep -n  "SharedMemory" Qflex/relax/env/vector/process1.py

# MS-Human-700 dimensions
python -c "import mujoco; m=mujoco.MjModel.from_xml_path(
  'Qflex/envs/ms_700_walk/MS-Human-700/MS-Human-700.xml');
  print(m.nq, m.nu, m.na, m.ntendon)"             # 85 700 700 700
\end{lstlisting}

\paragraph{Version pins at time of writing.}
MuscleMimic: \texttt{jax==0.7.2}, \texttt{flax==0.12.0}, \texttt{mujoco==3.4.0},
\texttt{mujoco-mjx==3.4.0}, \texttt{warp-lang==1.10.0}, \texttt{orbax-checkpoint==0.11.37}.
mjlab 1.4.0: \texttt{torch>=2.7.0}, \texttt{rsl-rl-lib==5.2.0},
\texttt{mujoco\textasciitilde{}=3.8.0}, \texttt{mujoco-warp>=3.8.0.3},
\texttt{warp-lang>=1.12.0}, \texttt{onnxscript>=0.5.4}.
Qflex (\texttt{relax} 0.1.0): \texttt{jax[cuda12]==0.4.27}, \texttt{dm-haiku},
\texttt{optax}, \texttt{gymnasium[mujoco]==1.0.0}, \texttt{dill==0.3.8}.

\paragraph{Caveats.} Qflex's \texttt{relax} contains twelve algorithms, of which
\texttt{flow\_exp.py} is Qflex proper and the remainder are reimplemented baselines;
statements about ``Qflex'' above refer to \texttt{FlowExp} unless noted. mjlab throughput
figures are not measured here and are quoted only as comparable in order of magnitude.
The MuscleMimic paper reports 123 joints / 72 DoF for MyoFullBody; the Qflex paper reports
206 joints for MS-Human-700 while MuJoCo loads 85 independent DoF --- both discrepancies
reflect anatomical joints being counted separately from independent coordinates.

\end{document}
